\section{研究背景与意义}

近年来，协作机器人、自主移动机器人与具身智能机器人在工业制造、仓储物流、医疗服务等领域加速落地，
应用形态也从传统生产线中的重复作业，
逐步扩展到人机协作装配、柔性分拣搬运、辅助诊疗与护理服务等更加开放的场景。
与固定工位自动化设备相比，
智能机器人需要在动态环境中持续感知外界变化、理解任务目标，
并将决策结果转化为稳定、可重复的物理动作，
因此对算法模型、基础软件和边缘计算平台提出了更高要求。
面向“十五五”时期未来产业发展需求，
具身智能正被视为连接人工智能与物理世界的重要方向：
它不仅关注模型在数字空间中的识别与生成能力，
更强调智能体在真实环境中的感知、交互、学习与执行能力。
随着大模型、低成本机器人硬件和开源机器人软件栈的发展，
如何将先进智能能力可靠部署到实际机器人系统中，
正在成为机器人产业化和规模化应用的重要基础问题。

在技术层面，视觉—语言—动作（Vision-Language-Action, VLA）大模型将视觉感知、语言理解与动作决策统一到端到端神经网络中，
使机器人能够更好地理解开放环境、人类指令和复杂操作任务。
这类模型进一步增强了机器人系统的泛化能力，
也推动机器人从封闭场景中的固定流程执行，走向更开放、更灵活的智能操作。

然而，VLA 模型和具身智能软件通常具有较高的计算、存储与实时性需求。
当其部署到树莓派等资源受限的边缘平台时，
推理延迟、传感器观测、动作执行和系统调度等环节都可能成为性能瓶颈。
因此，如何在边缘平台上高效、可靠地运行机器人智能软件，
已成为制约机器人大规模部署的关键问题。

从技术栈角度看，机器人系统通常由三个层次组成。
硬件层涵盖上位机（嵌入式计算板）、传感器（摄像头、力传感器）、执行器（电机、舵机）与机械结构，
决定机器人的物理感知与运动能力；
算法层包括视觉感知、策略规划与运动控制，
以 ACT\cite{act} 和 Diffusion Policy\cite{diffusionpolicy} 等端到端模型为代表；
基础软件层则由操作系统、通信中间件和推理框架构成，
向下管理硬件资源，向上支撑算法模型的运行。
本文的研究聚焦于基础软件层：操作系统的内核调度策略、系统调用路径与设备驱动模型，
以及推理框架的线程管理与算子调度，
共同决定着机器人系统的硬件生态兼容性、推理性能、决策准确性、运行可靠性和系统安全性，
是整个机器人软件栈性能与可信度的基础。

因此，对边缘机器人基础软件开展系统性性能分析具有重要研究意义。
一方面，它能够解释端到端控制循环中观测、推理、执行和等待等阶段的真实开销来源；
另一方面，它能够为推理运行时优化、操作系统调度改进和硬件 I/O 路径设计提供依据。
对于缺少 GPU、计算资源和调度余量有限的低成本机器人平台而言，
这类跨层次分析也是推动具身智能软件从实验室验证走向稳定部署的重要基础。
